\documentclass[10pt,twocolumn,a4paper]{article}
\usepackage[a4paper,margin=1.65cm,columnsep=0.55cm]{geometry}
\usepackage[T1]{fontenc}
\usepackage{newtxtext,newtxmath}
\usepackage{graphicx,booktabs,microtype,xcolor,amsmath,array}
\usepackage[hidelinks]{hyperref}
\usepackage[font=small,labelfont=bf,skip=4pt]{caption}
\usepackage{enumitem}
\setlist{nosep,leftmargin=*}
\raggedbottom
\setlength{\textfloatsep}{9pt plus 2pt minus 2pt}
\setlength{\floatsep}{7pt plus 2pt minus 2pt}
\setlength{\dbltextfloatsep}{9pt plus 2pt minus 2pt}
\graphicspath{{../../outputs/report001/figures/}}
\input{../../outputs/report001/numbers.tex}

% One accent colour for structure; the figure palette for the model names.
\definecolor{accent}{HTML}{0B3C6E}
\definecolor{released}{HTML}{D55E00}
\definecolor{ftall}{HTML}{0072B2}
\definecolor{confirmed}{HTML}{1A7A45}
\definecolor{exploratory}{HTML}{9A6A00}
\definecolor{descriptive}{HTML}{4F6D8A}
\definecolor{nullres}{HTML}{6B6963}
\newcommand{\status}[2]{\textcolor{#1}{\bfseries\scriptsize #2}}
\newcommand{\conf}{\status{confirmed}{CONFIRMED}}
\newcommand{\expl}{\status{exploratory}{EXPLORATORY}}
\newcommand{\desc}{\status{descriptive}{DESCRIPTIVE}}
\newcommand{\notest}{\status{nullres}{NOT ESTABLISHED}}
\newcommand{\R}{\ensuremath{R}}

\makeatletter
\renewcommand\section{\@startsection{section}{1}{\z@}%
  {-1.7ex plus -0.4ex minus -0.2ex}{0.7ex plus 0.2ex}%
  {\normalfont\large\bfseries\color{accent}}}
\renewcommand\subsection{\@startsection{subsection}{2}{\z@}%
  {-1.2ex plus -0.3ex minus -0.2ex}{0.4ex plus 0.1ex}%
  {\normalfont\normalsize\bfseries\color{accent}}}
\renewcommand\paragraph{\@startsection{paragraph}{4}{\z@}%
  {0.9ex plus 0.3ex minus 0.2ex}{-0.6em}{\normalfont\normalsize\bfseries}}
\makeatother

\title{\vspace{-1.6em}\textbf{Large Commands, Small Effects}\\[0.2em]
{\large A JEPA world model under-predicts block motion under large commands,\\[-0.15em]
and targeted fine-tuning largely repairs it\par}\vspace{-0.6em}}
\author{Omer Torres \qquad Inbal Solway\\[-0.1em]
{\small Deep Learning course project, 2026 \quad
\url{https://github.com/OmerTorres/jepa-action-effects}}\vspace{-0.6em}}
\date{}

\begin{document}
\twocolumn[{%
\renewcommand\twocolumn[1][]{#1}%
\maketitle
\begin{center}
\includegraphics[width=\textwidth]{teaser.pdf}
\captionof{figure}{\textbf{Large movement commands expose under-predicted block motion.}
``Push strength'' $m$ is command magnitude, not force.
(a)~One designed family, selected by median true motion rather than prediction error.
Left: real input frame; green T: goal marker. Right: block displacement along its true
direction. True motion (black) barely changes from 0.60 to 1.00, but LeWM predicts less,
even backwards at 1.00. Circles: readings from real end frames. This example's error is
more extreme than average.
(b)~Gain compares predicted with actual motion along the true direction (1: matching;
0.5: half). Across 300 fresh families: block (vermilion; solid: designed, dashed: demo),
agent (purple), real-frame readout (grey); bands: 95\% intervals. Strip: training-command
frequency; shaded: beyond 99.9\% of demo clips.
(c)~Mean block error at 0.60 on 240 fresh families (\S\ref{sec:finetune}); hatched:
zero training-loss weight at 0.60; dotted: released model.}
\label{fig:teaser}
\end{center}
\vspace{0.3em}
}]

\noindent\fcolorbox{accent!40}{accent!5}{\parbox{\dimexpr\columnwidth-2\fboxsep-2\fboxrule\relax}{%
\small\setlength{\baselineskip}{10pt}%
\textbf{\color{accent}Key terms}\par\smallskip
\begin{description}[nosep,leftmargin=0.6em,labelindent=0pt,font=\normalfont\bfseries]
\item[Latent / readout] an image feature vector / a network that maps it to positions.
\item[Released / FT] the unchanged published checkpoint / a fine-tuned copy;
an \emph{arm} is one training condition.
\item[Family / branch] one start with several pushes / one push; statistics group by family.
\item[Designed / demo] a constructed contact start / a start from an expert training demo.
\item[Real-frame check] readout accuracy on real images, required before interpreting
predictions. \textbf{Fresh, sealed}: new starts, test rules fixed before evaluation.
\end{description}}}

\begin{figure*}[t]
\centering
\includegraphics[width=\textwidth]{method.pdf}
\caption{\textbf{How we measure imagined physics.}
(a)~One test family, also used in Fig.~\ref{fig:ft}b. Three past frames, 0.5\,s apart,
are the input. Four one-second pushes start identically; real end frames share one zoom
(dashed: initial block outline). Agent travel increases with command magnitude, but
block displacement barely changes from 0.45 to 0.60.
(b)~Two predictor steps give an end latent. The block readout minus its start reading
gives $p$; real end frames give $r$; simulator motion gives $t$ (Eq.~\ref{eq:displacements}).
Real future frames are training targets only, never predictor inputs.
(c)~Gain $b$ (\S\ref{sec:readout}).}
\label{fig:method}
\end{figure*}

\begin{abstract}
\noindent Can a learned world model predict what happens when a robot receives movement
commands larger than those common in its training data? We test LeWorldModel (LeWM) in
PushT, a simulator where a disk pushes a T-shaped block. Starting repeatedly from the
same state, we vary the magnitude of a movement command held for one second and compare
the model's predicted block displacement with the simulator's actual outcome.
LeWM predicts feature vectors rather than positions, so a separately trained network
converts these vectors into positions; we check that network on real images too.
As command magnitude increases, decoded predictions increasingly understate block motion,
including designed scenes with nearly unchanged true displacement. The trend
recurs on fresh test sets and across five readouts. Fine-tuning on targeted
pushes reduces mean block-position error from \tdMeanFrozen{} to \tdMeanBalanced{} units
at command magnitude 0.60; equally long training on ordinary demonstrations shows no
established improvement. Most of the improvement remains when 0.60 is excluded from the
training loss. A short-horizon planning test does not establish a decision penalty.
\end{abstract}

\section{Introduction}
A world model predicts action outcomes to help a planner choose commands.
We ask whether its predictions correctly describe how a robot moves a block.

We test LeWM~\cite{maes2026lewm}, a joint-embedding predictive architecture (JEPA).
It encodes images into numerical feature vectors, called \emph{latents}, and predicts
future latents instead of images~\cite{assran2023ijepa,assran2025vjepa2}.
Like DINO-WM~\cite{zhou2024dinowm}, it can plan by comparing a predicted latent with
one encoded from a goal image. Good scores in this representation do not necessarily
guarantee accurate physical predictions~\cite{lambert2020objective}. We therefore train a
\emph{readout}: a network converting latents into positions, checked on real images
with simulator-supplied positions.

Our intervention changes the size of the robot's movement command, called \emph{push
strength} in the figures. This is not a force applied to the block. With the same start
state and commanded direction, a larger command changes how the robot approaches and
passes the block; contact geometry and contact time also affect the resulting motion.
We compare each prediction with its own simulator outcome, without assuming that larger
commands must cause larger block displacements. We hold each command for one second;
\S\ref{sec:task} explains this choice and the model's separate timing constraints.

\textbf{Q1:} does LeWM predict block motion accurately as commands become larger and
less common in its training demos? \textbf{Q2:} does an apparent shortfall survive
changes to the readout? \textbf{Q3:} can fine-tuning repair it, and with which data?
We summarise motion with a \emph{gain}: for example, predicting 6 units to the right
when the block really moves 10 gives gain 0.6. The aggregate version measures motion
along the true direction, not sideways errors (\S\ref{sec:readout}); it is not a measure
of improvement from training. Our training includes the readout and twelve fine-tuning
runs of LeWM's \ftParameters-parameter predictor.

\section{Background and related work}
LeWM trains its image encoder and action-conditioned predictor end to end from
pixels~\cite{maes2026lewm}. The Intervention Gap~\cite{vakalis2026gap} also audits LeWM's
response to actions, in DMC Cheetah. AD-WM~\cite{qiu2026adwm} and
Do-JEPA~\cite{resani2026dojepa} instead change training objectives to improve action
sensitivity. Kang et al.~\cite{kang2023extrapolate} report that networks can approach
constant predictions far outside their training inputs, motivating one explanation we
test. We follow PhysWeep~\cite{khanbayov2026physweep} in comparing explanations on
held-out data. Most are recent preprints; we have not reproduced their results.
Our contribution combines controlled pushes from identical starts, physical measurement
checks and targeted fine-tuning of the released PushT model.

\begin{table*}[t]
\centering\footnotesize
\setlength{\tabcolsep}{3.2pt}
\renewcommand{\arraystretch}{1.12}
\begin{tabular}{@{}>{\raggedright}p{1.7cm}>{\raggedright}p{2.95cm}>{\raggedright}p{2.8cm}>{\raggedright}p{3.85cm}>{\raggedright}p{2.95cm}l@{}}
\toprule
Test & Question it answers & Fresh test set & Check on the readout, before scoring & Answer & Status\tabularnewline
\midrule
E1 Strong push & Does LeWM get a strong push right? & 100 designed families, push 0.60 & readout error below 25\% of the true effect & No: gain \physSlopePredPt{} vs \physSlopeRealPt{} on real frames & \conf\tabularnewline
E2 Strength sweep & Is it the push strength or our unusual scenes? & 240 families: designed 0.30--0.60; demo commands $\times$1--$\times$2 & Real images recover 85--115\% of true motion (strengths pooled); $\geq$30 moving families per group; baseline checks* & Larger commands worsen motion underestimation in both scene types & \conf\tabularnewline
E3 Beyond range & What happens past the demos' range? & 300 families, six strengths up to 1.00 & At each strength 0.60--1.00, per scene: real-image error $<$25\% of true motion; 85--115\% motion recovered; $\geq$50 moving families & Motion recovery keeps falling ($R=\supSlopeArchiveHundredPt{}$ at 1.00); no explanation wins & \conf\tabularnewline
E4 Readouts & Could it be our measuring instrument? & re-decoding of E1--E2 & as in E1--E2 & All five readouts agree; shared bias remains possible & \expl\tabularnewline
E5 Agent vs block & Do agent and block behave alike? & E3 families & real-frame agent gain 0.85--1.15 & No: agent gain \supAgentChallengeSixtyPt, block \supSlopeChallengeSixtyPt{} at 0.60 & \desc\tabularnewline
E6 Fine-tuning & Can training fix it, and with which data? & 240 families + 120 ordinary clips & as E3, at 0.45 and 0.60 & Push families fix most; ordinary data does not & \conf\tabularnewline
E7 Planning & Does it change decisions? & 148 families & --- & No cost detected; test is coarse & \notest\tabularnewline
\bottomrule
\end{tabular}
\caption{\textbf{The study as a chain of questions.}
We check the readout on real images first. ``85--115\%'' requires the entire 95\% interval
for motion gain within 0.85--1.15 (1: correct motion along its true direction).
``Moving'': block displacement $>5.12$ units. $R$ is predicted-motion gain divided by
real-image gain (1: matching that reference; lower: less motion recovered).
E1 checks differences between pushes, so shared position offsets cancel.
*E2 also checks that strong designed pushes are under-predicted and unscaled recorded
commands are predicted faithfully. The four count groups are weak/strong designed
pushes and original/doubled recorded commands.
Status: \conf{}: prespecified test passed on fresh data; \expl{}: reanalysis;
\desc{}: not prespecified; \notest{}: prespecified effect not detected.}
\label{tab:studies}
\end{table*}

\section{Setup: task, data and measurement}
\nopagebreak[4]
\subsection{Task, model and training demos}\label{sec:task}
In PushT an agent disk of radius 15 pushes a T-shaped block (120 units across) on a
512$\times$512-unit table. Observations are 224$\times$224 RGB renders. Each command is a
two-dimensional vector in $[-1,1]^2$ that sets a target $100\times$ the command away from
the agent, which a proportional--derivative (PD) controller tracks. We use
the released checkpoint (\texttt{quentinll/lewm-pusht}: ViT-tiny encoder, 192-dimensional
latent, six-layer predictor with adaLN action conditioning, about 15M parameters) and the
pinned stable-worldmodel PushT simulator~\cite{maes2026swm}.

\paragraph{Why five commands, and why one second?} The simulator accepts a command every
0.1\,s. LeWM was trained with five consecutive commands per prediction step (0.5\,s),
and three input images spaced 0.5\,s apart; these are inherited interface choices.
We chose a two-prediction-step horizon: ten future commands, or one second, feeding the
first prediction back without future images. In a simulator-only pilot, 88 of 120 start
states produced block endpoints more than 5.12 units apart under different commands.
We retained this horizon as sufficient to reveal action-dependent motion;
we did not test its optimality.

\emph{Push strength} is command magnitude $m=\sqrt{x^2+y^2}$, not force, speed or
distance. Constant-command tests repeat the vector ten times. Each component lies
in $[-1,1]$, so length 1 is allowed in every direction (diagonals reach $\sqrt2$).
In Fig.~\ref{fig:method}a, raising $m$ from 0.45 to 0.60 increases agent travel
from 178 to 238 units, but the block moves about 24 in both cases: measurements,
not a distance formula.
The training demos comprise \demoEpisodes{} episodes and
\demoFrames{} frames. Over \demoWindows{} sampled one-second clips, the mean command has
median \demoCommandMedian, 90th percentile \demoCommandNinety{} and 99th percentile
\demoCommandNinetynine; only \supMapAbove{} clips exceed \supMapPninetynine, and none exceeds
\supMapMax. Our test pushes span 0.15 to 1.00, so 0.80 and 1.00 lie at and beyond the edge
of what the demos contain (mean command is not contact speed, so this does not prove such
interactions never occurred in training).

\subsection{Counterfactual families and fresh test sets}\label{sec:families}
A \emph{family} branches from one start into alternative pushes. Each branch resets
with the same random seed (fixing random choices), then replays the same earlier
commands (the \emph{warm-up}). Replay checks found identical starts. Agent--block
contacts are recorded every 0.01\,s: ten \emph{physics substeps} per 0.1\,s command.

\emph{Designed} scenes place the agent beside the block, offset sideways, and command
it toward the block. \emph{Demo} starts come from demonstrations with a nearly stationary
block. We either repeat a command in the demonstrator's average next direction at
different magnitudes, or multiply the next ten recorded commands by $1$, $1.5$ or $2$,
preserving their direction sequence. A branch is \emph{moved} if its final block
position is over 5.12 units from its start (1\% of table width).

Confirmatory hypotheses, metrics and decision rules were hash-sealed before collecting
fresh starts disjoint from earlier studies. Estimates carry 95\% intervals from 2,000
family bootstrap resamples, stratified by scene and paired across models; branches
share a start and are not independent. Figure examples follow a rule fixed before
plotting: true displacement closest to the cohort median.

\subsection{Reading physics out of latents}\label{sec:readout}
The readout $g$ is a PyTorch multilayer perceptron (Fig.~\ref{fig:readout}). With
$\rho=\mathrm{ReLU}$ and standardised input $\tilde z=(z-\mu_z)\oslash\sigma_z$,
\begin{equation}\label{eq:readout}
\begin{aligned}
h_1&=\rho(W_1\tilde z+b_1),\quad h_2=\rho(W_2h_1+b_2),\\
g(z)&=\mu_y+\sigma_y\odot(W_3h_2+b_3).
\end{aligned}
\end{equation}
Here $\odot,\oslash$ act elementwise; means and scales are fitted on training frames
only. The six outputs, in order, are block $(x,y)$, $\sin\phi,\cos\phi$ for block angle
$\phi$, and agent $(x,y)$. Let $g_{\mathrm{B}}$ select the first two outputs.
Adam ($10^{-3}$, batch 512, standardised targets, early-stop patience 20, at most 300 epochs)
fits \readoutTrainFrames{} uniformly sampled pose renders. We select among three MLP
seeds and a linear readout by validation block error: \readoutBlockError{} vs
\readoutRidgeError{} units on \readoutValFrames{} frames.

\begin{figure}[t]
\centering
\includegraphics[width=\columnwidth]{readout.pdf}
\caption{\textbf{Readout architecture} (Eq.~\ref{eq:readout}). Two 256-unit hidden layers
use ReLU; the final layer is linear. Standardisation and its inverse use training-set
statistics. The same frozen readout measures real and predicted latents.}
\label{fig:readout}
\end{figure}

For branch $i$, let $z_{0i}$ be the start latent, $\hat z_i$ the predicted end latent,
$o_i$ the real end image, and $q_{0i},q_i$ the simulator's start/end block positions.
We compare imagined, real-frame and true displacements:
\begin{equation}\label{eq:displacements}
\begin{aligned}
p_i&=g_{\mathrm B}(\hat z_i)-g_{\mathrm B}(z_{0i}),\\
r_i&=g_{\mathrm B}(E(o_i))-g_{\mathrm B}(z_{0i}),\quad t_i=q_i-q_{0i}.
\end{aligned}
\end{equation}
Eq.~\ref{eq:gains} pools the scored branches $i\in\mathcal I$:
\begin{equation}\label{eq:gains}
\begin{aligned}
b&=\frac{\sum_i p_i^\top t_i}{\sum_i\lVert t_i\rVert^2},\quad
b_{\mathrm{real}}=\frac{\sum_i r_i^\top t_i}{\sum_i\lVert t_i\rVert^2},\\
R&=b/b_{\mathrm{real}},\quad
\eta=\frac{\sum_i\lVert p_i\rVert^2}{\sum_i\lVert t_i\rVert^2}.
\end{aligned}
\end{equation}
Gain $b$ is a least-squares slope through the origin (1: matching motion along the true
direction); $R$ is a ratio of gains. Pooling avoids dividing by each small move; scores
require positive denominators. $R$ normalises by real-frame gain, but imagined-latent
readout bias can remain. $\mathcal I$ contains moved branches, except in the six-strength
study, which uses all branches. Gains ignore sideways error; alignment tracks it.
For endpoint accuracy we use
\begin{equation}\label{eq:endpoint}
e_i=\lVert g_{\mathrm B}(\hat z_i)-q_i\rVert_2,
\end{equation}
including unmoved branches. Fine-tuning results average training seeds and weight
families within each scene equally, then scenes equally. Real-frame readings provide
the reference through the same readout.


\section{What goes wrong}\label{sec:results}
\paragraph{E1: a strong push is under-predicted.} On 100 fresh designed families at
strength 0.60, predicted block gain is \physSlopePred, against
\physSlopeReal{} through the same readout on real end frames.

\paragraph{E2: the trend occurs in both scene types.} Across 240 new families,
designed-scene \R{} falls from \shiftRchThirty{} at 0.30 to
\shiftRchSixty{} at 0.60; when the demonstrators' own commands are scaled from $\times1$ to
$\times2$, it falls from \shiftRarOne{} to \shiftRarTwo{} (Table~\ref{tab:studies}).
Holding commands constant rather than varying them made no established difference
(\shiftdSar), within the prespecified equivalence band.

\paragraph{E3: the shortfall grows beyond the training range.} On 300 new families
pushed at six strengths up to 1.00 (Fig.~\ref{fig:teaser}b), the block gain falls from
\supSlopeChallengeSixty{} at 0.60 to \supSlopeChallengeHundred{} at 1.00 in designed scenes
and from \supSlopeArchiveSixty{} to \supSlopeArchiveHundred{} in demo scenes.
Median true displacement changes little (\supTrueChallengeSixty{} vs
\supTrueChallengeHundred{} units in designed scenes, \supTrueArchiveSixty{} vs
\supTrueArchiveHundred{} in demo scenes). Designed endpoints also barely change;
some demo families move further despite their stable median. Alignment changes
little until the strongest push, indicating mainly magnitude error. Three preregistered
explanations, fitted up to 0.60 and tested above, compare proportional shrinkage,
drift toward typical demo motion, and saturation at the strongest demo command. None fitted clearly
best: saturation beat proportional shrinkage (difference in squared error \supCminusA) but
not drift (\supCminusB); these differences are in squared table units.

\begin{figure}[t]
\centering
\includegraphics[width=\columnwidth]{robustness.pdf}
\caption{\textbf{Five different readouts see the same trend.} Change from weak to strong
pushes in energy ratio $\eta$ (Eq.~\ref{eq:gains}; squared, direction-free magnitude)
when the saved predictions of E1--E2 are re-decoded by
five readouts that extrapolate differently. All intervals lie below zero.}
\label{fig:robust}
\end{figure}

\paragraph{E4: the trend survives alternative readouts.} The readout recovers real motion at
every strength, even at 1.00 (real-frame gain \supRealChallengeHundred{} in designed and
\supRealArchiveHundred{} in demo scenes). Five readouts (three MLP seeds, linear ridge,
and non-extrapolating nearest neighbours) give the same downward trend
(Fig.~\ref{fig:robust}). Agreement supports the trend but cannot exclude shared
encoder/frame-bank bias. A readout-free method (render-and-encode pose inversion)
failed engineering checks; restricting predictions to those resembling real latents
left one readout inconclusive
(\robustOnChallengeRidge).

\paragraph{E5: agent and block behave differently.} Decoding the \emph{agent} from the
same imagined latents gives agent gains of \supAgentChallengeSixty{} (designed) and
\supAgentArchiveSixty{} (demo) at strength 0.60, and \supAgentChallengeEighty{} at 0.80,
while block gains are \supSlopeChallengeSixtyPt{} and \supSlopeArchiveSixtyPt. This
agent--block difference suggests interaction error, but aggregate agent gain cannot
rule out command-processing problems. We report agent gains only where real-frame
checks pass; at 1.00 most agents leave the table (only
\inViewChallengeHundred{} of designed and \inViewArchiveHundred{} of demo agents stay in
view). Extending this rule to the agent after seeing values makes E5 descriptive.

\section{Can training fix it?}\label{sec:finetune}
If rare strong interactions explain compression, targeted training should help more
than ordinary-data training, and nearby strengths might suffice. We test these predictions
by fine-tuning.

\paragraph{Setup.} We train the predictor $P_\theta$, the action encoder $A_\theta$ and the
prediction projector (\ftParameters{} parameters). The encoder, its projector, all
BatchNorm statistics and the readouts stay frozen. With cached real history latents and
action-block embeddings $\alpha_k=A_\theta(a_k)$ for two past and two candidate command
blocks (normalised and flattened), the rollout is
\begin{equation}\label{eq:rollout}
\begin{aligned}
\hat z_1 &= P_\theta\big([z_{-2},z_{-1},z_0],[\alpha_{-2},\alpha_{-1},\alpha_0]\big),\\
\hat z_2 &= P_\theta\big([z_{-1},z_0,\hat z_1],[\alpha_{-1},\alpha_0,\alpha_1]\big).
\end{aligned}
\end{equation}
$P_\theta$ here includes the prediction projector; gradients flow through $\hat z_1$.
For family $f$ and magnitude $m$, the branch loss and training objective are
\begin{equation}\label{eq:loss}
\begin{aligned}
\ell_{fm}&=\frac{1}{2\cdot192}\sum_{k=1}^2\lVert\hat z_{fmk}-E(o_{fmk})\rVert_2^2,\\
\mathcal L&=\frac{1}{2}\sum_{s=1}^2\frac{1}{|\mathcal F_s|}
\sum_{f\in\mathcal F_s}\sum_m w_m\ell_{fm}.
\end{aligned}
\end{equation}
Here $\mathcal F_s$ is the batch's families in scene $s$ and $\sum_m w_m=1$;
the weights define each training condition (\emph{arm}, Table~\ref{tab:ft}). AdamW (learning
rate $10^{-5}$, selected from $\{10^{-5},3{\cdot}10^{-5},10^{-4}\}$ by validation latent
error; weight decay 0.01; clipping 1.0) runs 2,000 updates, each with 16 families
$\times$ 4 strengths. Push-family arms pair batches by seed (three seeds). Training uses 600 families
(300 per scene), strengths 0.15--0.60, and 120 validation families. All-strength and
demo-weighted runs are reused from the earlier study. The ordinary arm averages loss
over 64 recorded demo clips per update, from 2,400 clips. Testing uses 240 fresh
families and 120 ordinary clips; all four readout checks passed. The omitted arm uses
$w=(1/3,1/3,1/3,0)$ in ascending magnitude order: Eq.~\ref{eq:loss} does not divide by
four again, leaving 48 nonzero-weight branches out of 64 per update. Strength
0.60 still informed family admission and the earlier learning-rate
choice, and appears in validation. Every arm uses its fixed final checkpoint.
Thus 0.60 is excluded from updates, not all development.

\begin{table}[t]
\centering\footnotesize
\setlength{\tabcolsep}{3pt}
\begin{tabular}{@{}l>{\raggedright\arraybackslash}p{2.25cm}rrrr@{}}
\toprule
& & \multicolumn{3}{c}{strength 0.60} & ordinary \\
\cmidrule(lr){3-5}
Model & fine-tuned on & error & \R{} des. & \R{} demo & error \\
\midrule
\textcolor{released}{released} & --- & \tdMeanFrozen & \tdRatioChallengeFrozen & \tdRatioArchiveFrozen & \tdOrdMeanFrozen \\
FT ordinary & 2,400 ordinary demo clips & \tdMeanArchOnly & \tdRatioChallengeArchOnly & \tdRatioArchiveArchOnly & \tdOrdMeanArchOnly \\
FT demo-wt. & push families, weighted like the demos & \tdMeanArchWeighted & \tdRatioChallengeArchWeighted & \tdRatioArchiveArchWeighted & \tdOrdMeanArchWeighted \\
FT w/o 0.60 & push families, 0.15--0.45 only & \tdMeanOmitted & \tdRatioChallengeOmitted & \tdRatioArchiveOmitted & \tdOrdMeanOmitted \\
\textcolor{ftall}{FT all} & push families, all four strengths & \tdMeanBalanced & \tdRatioChallengeBalanced & \tdRatioArchiveBalanced & \tdOrdMeanBalanced \\
\midrule
real frames & --- & & & & \tdOrdMeanReal \\
\bottomrule
\end{tabular}
\caption{\textbf{Fine-tuning arms and results} (fresh data, seeds averaged). Error: mean
block endpoint error (Eq.~\ref{eq:endpoint}) in table units (table: 512; block: 120), scenes weighted
equally. FT ordinary controls for training as such. \R: gain relative to real
frames (1 = as good as the instrument). Demo weights are the frequencies of the four
strengths in demo clips (.51/.34/.13/.02).}
\label{tab:ft}
\end{table}

\begin{figure*}[t]
\centering
\includegraphics[width=\textwidth]{finetuning.pdf}
\caption{\textbf{Push-family fine-tuning removes most compression, including at a strength
excluded from the loss; ordinary data does not.} Seed means.
(a)~\R{} on 240 fresh families; shaded: zero loss weight at 0.60 in ``w/o 0.60''.
(b)~Fig.~\ref{fig:method}a's family at 0.60.
(c)~Validation latent error during training: 120 families (60 designed, 60 demo).}
\label{fig:ft}
\end{figure*}

\paragraph{Results.} Fine-tuning on all four strengths lowers the mean block error at
strength 0.60 from \tdMeanFrozen{} to \tdMeanBalanced{} units (difference
\tdFrozenMinusBal) and raises \R{} from \tdRatioChallengeFrozen{} to \tdRatioChallengeBalanced{} in designed and
from \tdRatioArchiveFrozen{} to \tdRatioArchiveBalanced{} in demo scenes
(Table~\ref{tab:ft}, Fig.~\ref{fig:ft}a). Fine-tuning as long on ordinary demo clips gives
no established improvement (\tdFrozenMinusArchOnly); its interval rules out improvements above
\tdArchOnlyBound{} units, against \tdFrozenMinusBalPt{} for push families. Leaving
strength 0.60 out of the training loss leaves a small but preregistered deficit there
(\tdOmitMinusBal{} units) and no established deficit at the trained strength 0.45
(\tdOmitMinusBalInSupport); the held-out model still recovers \tdTransferPercent{} of the
full error reduction (\tdTransfer, descriptive). Its deficit does not separate coverage
from the number of nonzero-weight branches. Weighting
the strengths like the demos changed little (\tdMeanArchWeighted). The improvement has a
price: on ordinary demo clips the all-strength model is \tdOrdBalMinusFrozen{} units worse
than the released one (\tdOrdMeanBalanced{} vs \tdOrdMeanFrozen; real frames give
\tdOrdMeanReal). Its cause is not isolated; mixing ordinary clips into targeted training
is a candidate mitigation. Validation curves (Fig.~\ref{fig:ft}c) show stable optimisation. An
earlier fine-tuning study agrees (\trainBalMinusFrozen{} units) but is only descriptive,
because its readout check failed in demo scenes.

\section{Does it matter for decisions?}
We rank six candidates $c\in\mathcal C$ by predicted-end-to-goal latent distance.
For goal position $q^*$, let $d_c=\lVert q_c-q^*\rVert_2$ be candidate $c$'s true
block-to-goal distance. The chosen candidate $\hat c$ incurs
\begin{equation}\label{eq:regret}
\mathrm{regret}(\hat c)=d_{\hat c}-\min_{c\in\mathcal C}d_c.
\end{equation}
On 148 fresh families, real-latent choice beats random (regret difference
\planRealMinusRandom{} units; Eq.~\ref{eq:regret}). Imagined-minus-real regret is
\planJepaMinusReal: no positive penalty is established, but the interval allows costs
up to \planBound{} units. Each candidate's privileged goal image preserves its realised final
agent position and block angle, moving only the block to the goal. Predicted agent/angle
errors can still affect rankings, so this tests total latent prediction error, not block
error alone. Six candidates and a two-step horizon make the test coarse; the null does
not establish that compression is harmless.

\section{Discussion}
\paragraph{What we can conclude.} LeWM increasingly under-predicts block motion under
large commands while agent gain stays near one, even where designed-scene outcomes
barely change. Targeted fine-tuning removes most compression, including at a strength
excluded from the loss; equally long ordinary-data training does not. This favours
interaction coverage over an architectural limit, but neither proves pretraining-data
causation nor locates the repair's mechanism.

\paragraph{Insights.} Ordinary clips hide the shortfall: model error
(\tdOrdMeanFrozen{}) is near readout error (\tdOrdMeanReal). Controlled pushes expose
it; targeted training trades push accuracy for ordinary-data accuracy.
Three earlier readouts failed their real-frame checks.

\paragraph{Limitations and next steps.} One checkpoint and environment; a two-step horizon
and constant pushes; readouts that share one encoder and frame bank; some exploratory
analyses on opened data; fewer nonzero-weight examples in the omitted-strength arm; fine-tuning
seeds rather than independently pretrained models; a coarse planning test. Next: a
readout-free instrument, more environments and checkpoints, and pretraining without the
strongest commands.

\paragraph{Reproduction.} \texttt{reports/report001/REPRODUCE.md} gives commands to
rebuild all numbers, tables, figures and this PDF from tracked sources.

{\footnotesize
\begin{thebibliography}{99}\setlength{\itemsep}{0pt}
\bibitem{assran2023ijepa} M.~Assran et al. Self-supervised learning from images with a
joint-embedding predictive architecture. arXiv:2301.08243, 2023.
\bibitem{assran2025vjepa2} M.~Assran et al. V-JEPA 2: Self-supervised video models
enable understanding, prediction and planning. arXiv:2506.09985, 2025.
\bibitem{maes2026lewm} L.~Maes, Q.~Le~Lidec, D.~Scieur, Y.~LeCun, R.~Balestriero.
LeWorldModel: Stable end-to-end joint-embedding predictive architecture from pixels.
arXiv:2603.19312, 2026.
\bibitem{maes2026swm} L.~Maes et al. stable-worldmodel: A platform for reproducible world
modeling research and evaluation. arXiv:2605.21800, 2026.
\bibitem{zhou2024dinowm} G.~Zhou, H.~Pan, Y.~LeCun, L.~Pinto. DINO-WM: World models on
pre-trained visual features enable zero-shot planning. arXiv:2411.04983, 2024.
\bibitem{lambert2020objective} N.~Lambert, B.~Amos, O.~Yadan, R.~Calandra. Objective
mismatch in model-based reinforcement learning. arXiv:2002.04523, 2020.
\bibitem{vakalis2026gap} D.~Vakalis. The intervention gap in latent world models.
arXiv:2608.29998, 2026.
\bibitem{kang2023extrapolate} K.~Kang, A.~Setlur, C.~Tomlin, S.~Levine. Deep neural
networks tend to extrapolate predictably. ICLR 2024 (arXiv:2310.00873).
\bibitem{khanbayov2026physweep} R.~Khanbayov, H.~Kurban. PhysWeep: Does a video generator
realize the physics you ask for? arXiv:2609.06207, 2026.
\bibitem{resani2026dojepa} H.~Resani, J.~Q.~Shi. Do-JEPA: From masking to intervention in
latent world models. arXiv:2609.37378, 2026.
\bibitem{qiu2026adwm} J.~Qiu et al. AD-WM: Action-discriminative world models for
counterfactual model predictive control. arXiv:2609.30264, 2026.
\end{thebibliography}}
\end{document}
